% Custom macros common to all reports:
{{ producercommand }}


%%%%%%%%%%%%%%%%%%%%%%%%%%%%%%%%%%%%%%%%%
%% The document starts here
%%%%%%%%%%%%%%%%%%%%%%%%%%%%%%%%%%%%%%%%%

% Every unit of Zone 1 inputs its own copy of this report into the final one, so its macros are
% provided, defined by the first copy only.

% The version Type C/D compares the curves of InternalNode1, the output of the unit after its
% transformer, where the fiche observes the results of the Fiche I18, and the captions say so
\providecommand{\GridCurvesDescriptionNodeOne}{
    \figuredrawnzonefalse
    \GridCurveDescriptionZone{fig_V}{Voltage at InternalNode1.}
    \GridCurveDescriptionZone{fig_P}{Active power at InternalNode1.}
    \GridCurveDescriptionZone{fig_Q}{Reactive power at InternalNode1.}
    \GridCurveDescriptionZone{fig_Ip}{Active current at InternalNode1.}
    \GridCurveDescriptionZone{fig_Iq}{Reactive current at InternalNode1.}
    \GridCurveDescriptionZone{fig_WRef}{Frequency of the converter.}
    \iffiguredrawnzone\else\item No figure: the test produced no curve.\fi
}

% #1 = what the line Zs does at the event, #2 = the schematic of its state before the event
\providecommand{\GridCircuitZoneSwitchedLine}[2]{
    The grid model and its operational point is as in the following schematic, where the line
    Zs#1:
    \begin{center}
        \includestandalone[width=0.8\textwidth]{#2}
    \end{center}
    \begin{center}
        \small \textbf{Note: This schematic is only a reminder of the test setup on the TSO's
        side --- the Producer's side may vary, depending on the user-provided model.}
    \end{center}
}

\providecommand{\GridCurvesNodeOne}{%
    \setcounter{gridfigure}{0}%
    \noindent
    \GridCurveZone{fig_V}{Voltage magnitude measured at InternalNode1.}
    \GridCurveZone{fig_P}{Active power P at InternalNode1.}
    \GridCurveZone{fig_Q}{Reactive power Q at InternalNode1.}
    \GridCurveZone{fig_Ip}{Active current Ip at InternalNode1.}
    \GridCurveZone{fig_Iq}{Reactive current Iq at InternalNode1.}
    \GridCurveZone{fig_WRef}{Frequency of the converter.}
}

    \section{Results for PCS RTE-I16 Zone 1, version Type C/D ({{producer}})}

    In this report, \textbf{InternalNode1} denotes the internal connection point of the
    Zone 1 unit model inside the aggregated plant model (the node named \emph{Node1} in
    the DTR). The name \emph{PDR} is reserved for the connection point of the complete
    installation to RTE's grid. The Fiche I16.2 observes the references of Zone 1, the results
    of the Fiche I18, at the output of the unit after its transformer, so every magnitude this
    report compares (and the corresponding curve signal names) refers to InternalNode1.

    \BLOCK{for subReport in subReports}
    {{subReport}}
    \newpage
    \BLOCK{endfor}

